\documentclass[UTF8]{ctexart}
\usepackage{xeCJK}
\usepackage{geometry}
\geometry{left=2cm,right=2cm,top=2.5cm,bottom=2.5cm}
\setlength{\headheight}{12.7pt}

\usepackage{titlesec}
\titleformat{\section}
  {\normalfont\large\bfseries}
  {\thesection}
  {1em}
  {}
\titlespacing*{\section}
  {0pt}
  {3.5ex plus 1ex minus .2ex}
  {2.3ex plus .2ex}

\usepackage{amsmath}
\usepackage{amssymb}
\usepackage{amsthm}
\usepackage{enumitem}
\setlist[enumerate]{itemsep=0pt,topsep=0pt,parsep=0pt,leftmargin=0pt,itemindent=2.5em}
\setlist[itemize]{itemsep=0pt,topsep=0pt,parsep=0pt,leftmargin=0pt,itemindent=2.5em}
\usepackage{fancyhdr}
\pagestyle{fancy}
\fancyhf{}
\usepackage{graphicx}
\usepackage{hyperref}
\usepackage{indentfirst}
\usepackage{lmodern}


\begin{document}
\setlength{\abovedisplayskip}{1pt}
\setlength{\belowdisplayskip}{1pt}

\section{部分正确性注释程序}

\hypertarget{sec:定义1.1}{}
\noindent\textbf{定义1.1 (部分正确性注释程序)}

对断言$P,Q$, 程序$S$和某字符串$S^*$, 递归定义
\[
    \{P\}\,S^*\,\{Q\}\ \vdash_{\mathrm{PO}}\ \{P\}\,S\,\{Q\},
\]
称$\{P\}\,S^*\,\{Q\}$是$\{P\}\,S\,\{Q\}$的\textbf{部分正确性注释程序},
$\{P\},\{Q\}$分别为前后断言($\mathrm{pre},\mathrm{post}$),
$S$为骨架程序($\mathrm{sk}$):
\vspace{1pt}
\begin{enumerate}
    \item 当$P\to Q$时, $\{P\}\,\mathrm{skip}\,\{Q\}\vdash_{\mathrm{PO}}
    \{P\}\,\mathrm{skip}\,\{Q\}$,
    \item 当$P\to Q[u:=t]$时, $\{P\}\,u:=t\,\{Q\}\vdash_{\mathrm{PO}}
    \{P\}\,u:=t\,\{Q\}$,
    \item 如果$\{P\}\,S_1^*\,\{R\}\vdash_{\mathrm{PO}}\{P\}\,S_1\,\{R\}$且
    $\{R\}\,S_2^*\,\{Q\}\vdash_{\mathrm{PO}}\{R\}\,S_2\,\{Q\}$, 则
    \vspace{1pt}
    \[
        \{P\}\,S_1^*\,\{R\}\,S_2^*\,\{Q\}\vdash_{\mathrm{PO}}
        \{P\}\,S_1;S_2\,\{Q\},\\[1pt]
    \]
    \item 如果$\{P\land B\}\,S_1^*\,\{Q\}\vdash_{\mathrm{PO}}
    \{P\land B\}\,S_1\,\{Q\}$且$\{P\land\neg\,B\}\,S_2^*\,\{Q\}
    \vdash_{\mathrm{PO}}\{P\land\neg\,B\}\,S_2\,\{Q\}$, 则
    \vspace{1pt}
    \[\begin{aligned}
        &\,\,\{P\}\,\mathbf{if}\,\,B\,\,\mathbf{then}\,\,
        \{P\land B\}\,S_1^*\,\{Q\}\,\,\mathbf{else}\,\,\{P\land\neg\,B\}\,S_2^*
        \,\{Q\}\,\,\mathbf{fi}\,\,\{Q\}\\[-3pt]
        \,&\vdash_{\mathrm{PO}}\{P\}\,\mathbf{if}\,\,B\,\,\mathbf{then}\,\,S_1
        \,\,\mathbf{else}\,\,S_2\,\,\mathbf{fi}\,\{Q\},\\[2pt]
    \end{aligned}\]
    \item 如果$P\to I$, $\{I\land B\}\,S^*\,\{I\}
    \vdash_{\mathrm{PO}}\{I\land B\}\,S\,\{I\}$且$I\land\neg\,B\to Q$, 则
    \vspace{1pt}
    \[\begin{aligned}
        &\,\,\{P\}\{\mathbf{inv}:I\}\,\mathbf{while}\,\,B\,\,
        \mathbf{do}\,\,\{I\land B\}\,S^*\,\{I\}
        \,\,\mathbf{od}\,\{Q\}\,\\[-3pt]
        \,&\vdash_{\mathrm{PO}}\,\{P\}\,\mathbf{while}\,\,B
        \,\,\mathbf{do}\,\,S\,\,\mathbf{od}\,\{Q\}.
    \end{aligned}\]
\end{enumerate}

\vspace{15pt}

\hypertarget{sec:定理1.1}{}
\noindent\textbf{定理1.1 (正规形定理)}

对任意的断言$P,Q$和程序$S$, $\vdash_{\mathrm{PW}}\{P\}\,S\,\{Q\}$
当且仅当存在$\{P\}\,S^*\,\{Q\}\vdash_{\mathrm{PO}}\{P\}\,S\,\{Q\}$.

\vspace{15pt}

\hypertarget{sec:定义1.2}{}
\noindent\textbf{定义1.2 (部分正确性注释程序上的操作)}

把$\{P\}\,\mathrm{E}\,\{P\}$也视为自身的注释程序,
递归定义注释程序带状态的一步转移关系:
\begin{enumerate}
    \item $\langle\{P\}\,\mathrm{skip}\,\{Q\},\sigma\rangle
    \overset{\mathrm{def}}{\to}
    \langle\{Q\}\,\mathrm{E}\,\{Q\},\sigma\rangle$,
    \item $\langle\{P\}\,u:=t\,\{Q\},\sigma\rangle\overset{\mathrm{def}}{\to}
    \langle\{Q\}\,\mathrm{E}\,\{Q\},\sigma[u:=\sigma(t)]\rangle$,
    \vspace{3pt}
    \item 任给注释程序$\{P\}\,S_1^*\,\{R\}$和$\{R\}\,S_2^*\,\{Q\}$,
    如果$\langle\{P\}\,S_1^*\,\{R\},\sigma\rangle\to
    \langle\{P'\}\,S_1'\,^*\,\{R\},\sigma'\rangle$, 那么
    \vspace{1pt}
    \[
        \langle\{P\}\,S_1^*\{R\}\,S_2^*\,\{Q\},\sigma\rangle
        \overset{\mathrm{def}}{\to}
        \langle\{P'\}\,S_1'\,^*\,\{R\}\,S_2^*\,\{Q\},\sigma'\rangle,\\[1pt]
    \]
    \item 任给注释程序$\{P\}\,S_1^*\,\{R\}$和$\{R\}\,S_2^*\,\{Q\}$,
    如果$\langle\{P\}\,S_1^*\,\{R\},\sigma\rangle\to
    \langle\{R\}\,\mathrm{E}\,\{R\},\sigma'\rangle$, 那么
    \vspace{1pt}
    \[
        \langle\{P\}\,S_1^*\{R\}\,S_2^*\,\{Q\},\sigma\rangle
        \overset{\mathrm{def}}{\to}
        \langle\{R\}\,S_2^*\,\{Q\},\sigma'\rangle,\\[1pt]
    \]
    \item 任给注释程序$\{P\land B\}\,S_1^*\,\{Q\}$和
    $\{P\land\neg\,B\}\,S_2^*\,\{Q\}$, 如果$\sigma\models B$, 则
    \vspace{1pt}
    \[
        \,\,\langle\{P\}\,\mathbf{if}\,\,B\,\,\mathbf{then}\,\,
        \{P\land B\}\,S_1^*\,\{Q\}\,\,\mathbf{else}\,\,\{P\land\neg\,B\}\,S_2^*
        \,\{Q\}\,\,\mathbf{fi}\,\,\{Q\},\sigma\rangle\overset{\mathrm{def}}{\to}
        \langle\{P\land B\}\,S_1^*\,\{Q\},\sigma\rangle,\\[1pt]
    \]
    \item 任给注释程序$\{P\land B\}\,S_1^*\,\{Q\}$和
    $\{P\land\neg\,B\}\,S_2^*\,\{Q\}$, 如果$\sigma\not\models B$, 则
    \vspace{1pt}
    \[
        \,\,\langle\{P\}\,\mathbf{if}\,\,B\,\,\mathbf{then}\,\,
        \{P\land B\}\,S_1^*\,\{Q\}\,\,\mathbf{else}\,\,\{P\land\neg\,B\}\,S_2^*
        \,\{Q\}\,\,\mathbf{fi}\,\,\{Q\},\sigma\rangle\overset{\mathrm{def}}{\to}
        \langle\{P\land\neg\,B\}\,S_2^*\,\{Q\},\sigma\rangle,\\[1pt]
    \]
    \item 任给$P\to I$, 注释程序$\{I\land B\}\,S^*\,\{I\}$和
    $I\land\neg\,B\to Q$, 如果$\sigma\models B$, 则
    \vspace{1pt}
    \[\begin{aligned}
        &\langle\{P\}\{\mathbf{inv}:I\}\,\mathbf{while}\,\,B\,\,\mathbf{do}\,\,
        \{I\land B\}\,S^*\,\{I\}\,\,\mathbf{od}\,\{Q\},
        \sigma\rangle\overset{\mathrm{def}}{\to}\\[-3pt]
        &\langle\{I\land B\}\,S^*\,\{I\}\{\mathbf{inv}:I\}\,
        \mathbf{while}\,\,B\,\,\mathbf{do}\,\,
        \{I\land B\}\,S^*\,\{I\}\,\,\mathbf{od}\,\{Q\},
        \sigma\rangle,\\[1pt]
    \end{aligned}\]
    \item 任给$P\to I$, 注释程序$\{I\land B\}\,S^*\,\{I\}$和
    $I\land\neg\,B\to Q$, 如果$\sigma\not\models B$, 则
    \vspace{1pt}
    \[
        \langle\{P\}\{\mathbf{inv}:I\}\,\mathbf{while}\,\,B\,\,\mathbf{do}\,\,
        \{I\land B\}\,S^*\,\{I\}\,\,\mathbf{od}\,\{Q\},
        \sigma\rangle\overset{\mathrm{def}}{\to}
        \langle\{Q\}\,\mathrm{E}\,\{Q\},\sigma\rangle.
    \]
\end{enumerate}





\newpage

\hypertarget{sec:定理1.2}{}
\noindent\textbf{定理1.2}

对任意的$\langle G,\sigma\rangle\to\langle G',\sigma'\rangle$,

\begin{enumerate}
    \item $\mathrm{post}(G')=\mathrm{post}(G)$,
    \item 如果$\mathrm{sk}(G')=\mathrm{E}$,
    那么$G$一定形如$\{P\}\,\mathrm{E}\,\{P\}$,
    \item $\langle\mathrm{sk}(G),\sigma\rangle\to
    \langle\mathrm{sk}(G'),\sigma'\rangle$,
    \item 如果$\sigma\models\mathrm{pre}(G)$,
    那么$\sigma'\models\mathrm{pre}(G')$.
\end{enumerate}

\noindent{\kaishu 证明.}

\hspace{-0.71em} {\kaishu 1. 2.} 显然. \,\, {\kaishu 3. 4.} 对转移操作归纳证明.

\hspace{0.5em} i) $P\to Q$且$G=\{P\}\,\mathrm{skip}\,\{Q\}$,
则$G'=\{Q\}\,\mathrm{E}\,\{Q\}$, $\sigma'=\sigma$. 显然
\[
    \langle\mathrm{skip},\sigma\rangle\to
    \langle\mathrm{E},\sigma\rangle,
\]
并且如果$\sigma\models P$, 那么$\sigma\models Q$.

\hspace{0.5em} ii) 有$Q_1$使得$Q_1[u:=t]$是自由替换,
$P\to Q_1[u:=t]$且$Q_1\to Q$, $G=\{P\}\,u:=t\,\{Q\}$, 则
\[
    G'=\{Q\}\,\mathrm{E}\,\{Q\},\quad\sigma'=\sigma[u:=\sigma(t)],
\]

显然$\langle u:=t,\sigma\rangle\to
\langle\mathrm{E},\sigma[u:=\sigma(t)]\rangle$,
并且如果$\sigma\models P$, 那么$\sigma\models Q_1[u:=t]$, 从而
$\sigma'=\sigma[u:=t]\models Q_1$, 即$\sigma'\models Q$.

\vspace{1pt}
\hspace{0.5em} iii) $G=\{P\}\,S_1^*\,\{R\}\,S_2^*\{Q\}$,
假设$\langle\{P\}\,S_1^*\,\{R\},\sigma\rangle\to
\langle\{P'\}\,S_1'\,^*\,\{R\},\sigma'\rangle$并且
\[
    \langle S_1,\sigma\rangle\to\langle S_1',\sigma'\rangle,\quad
    \sigma\models P\to\sigma'\models P',
\]
那么$G'=\{P'\}\,S_1'\,^*\,\{R\}\,S_2^*\,\{Q\}$, 此时
$\langle S_1;S_2,\sigma\rangle\to\langle S_1';S_2,\sigma'\rangle$.

\hspace{0.5em} iv) $G=\{P\}\,S_1^*\,\{R\}\,S_2^*\{Q\}$,
假设$\langle\{P\}\,S_1^*\,\{R\},\sigma\rangle\to
\langle\{R\}\mathrm{E}\,\{R\},\sigma'\rangle$并且
\[
    \langle S_1,\sigma\rangle\to\langle\mathrm{E},\sigma'\rangle,\quad
    \sigma\models P\to\sigma'\models R,
\]
那么$G'=\{R\}\,S_2^*\,\{Q\}$, 此时
$\langle S_1;S_2,\sigma\rangle\to\langle S_2,\sigma'\rangle$.

\vspace{1pt}
\hspace{0.5em} v) $G=\{P\}\,\mathbf{if}\,\,B\,\,\mathbf{then}\,\,
\{P\land B\}\,S_1^*\,\{Q\}\,\,\mathbf{else}\,\,\{P\land\neg\,B\}\,S_2^*
\,\{Q\}\,\,\mathbf{fi}\,\,\{Q\}$且$\sigma\models B$, 则
\[
    G'=\{P\land B\}\,S_1^*\,\{Q\},\quad\sigma'=\sigma,
\]
显然$\langle\,\mathbf{if}\,\,B\,\,\mathbf{then}\,\,S_1\,\,\mathbf{else}\,\,
S_2\,\,\mathbf{fi},\sigma\rangle\to\langle S_1,\sigma\rangle$,
并且如果$\sigma\models P$, 那么$\sigma'=\sigma\models P\land B$.

\vspace{1pt}
\hspace{0.5em} vi) $G=\{P\}\,\mathbf{if}\,\,B\,\,\mathbf{then}\,\,
\{P\land B\}\,S_1^*\,\{Q\}\,\,\mathbf{else}\,\,\{P\land\neg\,B\}\,S_2^*
\,\{Q\}\,\,\mathbf{fi}\,\,\{Q\}$且$\sigma\not\models B$, 则
\[
    G'=\{P\land\neg\,B\}\,S_2^*\,\{Q\},\quad\sigma'=\sigma,
\]
显然$\langle\,\mathbf{if}\,\,B\,\,\mathbf{then}\,\,S_1\,\,\mathbf{else}\,\,
S_2\,\,\mathbf{fi},\sigma\rangle\to\langle S_2,\sigma\rangle$,
并且如果$\sigma\models P$, 那么$\sigma'=\sigma\models P\land\neg\,B$.

\vspace{1pt}
\hspace{0.5em} vii) $P\to I$, $I\land\neg\,B\to Q$,
$G=\{P\}\{\mathbf{inv}:I\}\,\mathbf{while}\,\,B\,\,
\mathbf{do}\,\,\{I\land B\}\,S^*\,\{I\}\,\,\mathbf{od}\,\{Q\}$
且$\sigma\models B$, 则
\[
    G'=\{I\land B\}\,S^*\,\{I\}\{\mathbf{inv}:I\}\,
    \mathbf{while}\,\,B\,\,\mathbf{do}\,\,
    \{I\land B\}\,S^*\,\{I\}\,\,\mathbf{od}\,\{Q\},\quad
    \sigma'=\sigma,
\]
显然
\vspace{-5pt}
\[
    \langle\,\mathbf{while}\,\,B\,\,\mathbf{do}\,\,S\,\,\mathbf{od},
    \sigma\rangle\to
    \langle S;\mathbf{while}\,\,B\,\,\mathbf{do}\,\,S\,\,\mathbf{od},
    \sigma\rangle,\\[-5pt]
\]
并且如果$\sigma\models P$, 那么$\sigma\models I$,
从而$\sigma'=\sigma\models I\land B$.

\vspace{1pt}
\hspace{0.5em} viii) $P\to I$, $I\land\neg\,B\to Q$,
$G=\{P\}\{\mathbf{inv}:I\}\,\mathbf{while}\,\,B\,\,
\mathbf{do}\,\,\{I\land B\}\,S^*\,\{I\}\,\,\mathbf{od}\,\{Q\}$
且$\sigma\not\models B$, 则
\[
    G'=\{Q\}\,\mathrm{E}\,\{Q\},\quad\sigma'=\sigma,
\]
显然
\vspace{-5pt}
\[
    \langle\,\mathbf{while}\,\,B\,\,\mathbf{do}\,\,S\,\,\mathbf{od},
    \sigma\rangle\to
    \langle\mathrm{E},\sigma\rangle,\\[-5pt]
\]
并且如果$\sigma\models P$, 那么$\sigma\models I$,
从而$\sigma'=\sigma\models I\land\neg\,B$, 即$\sigma'\models Q$.
\hfill $\qedsymbol$





\newpage

\section{完全正确性注释程序}

\hypertarget{sec:定义2.1}{}
\noindent\textbf{定义2.1 (完全正确性注释程序)}

对断言$P,Q$, 程序$S$和某字符串$S^*$, 递归定义
\[
    \{P\}\,S^*\,\{Q\}\ \vdash_{\mathrm{TO}}\ [P]\,S\,[Q],
\]
称$\{P\}\,S^*\,\{Q\}$是$[P]\,S\,[Q]$的\textbf{完全正确性注释程序},
前后断言与骨架程序的记法$\mathrm{pre},\mathrm{post},\mathrm{sk}$同上节:
\vspace{1pt}
\begin{enumerate}
    \item 当$P\to Q$时, $\{P\}\,\mathrm{skip}\,\{Q\}\vdash_{\mathrm{TO}}
    [P]\,\mathrm{skip}\,[Q]$,
    \item 当$P\to Q[u:=t]$时, $\{P\}\,u:=t\,\{Q\}\vdash_{\mathrm{TO}}
    \{P\}\,u:=t\,\{Q\}$,
    \item 如果$\{P\}\,S_1^*\,\{R\}\vdash_{\mathrm{TO}}[P]\,S_1\,[R]$且
    $\{R\}\,S_2^*\,\{Q\}\vdash_{\mathrm{TO}}[R]\,S_2\,[Q]$, 则
    \vspace{1pt}
    \[
        \{P\}\,S_1^*\,\{R\}\,S_2^*\,\{Q\}\vdash_{\mathrm{TO}}
        [P]\,S_1;S_2\,[Q],\\[1pt]
    \]
    \item 如果$\{P\land B\}\,S_1^*\,\{Q\}\vdash_{\mathrm{TO}}
    [P\land B]\,S_1\,[Q]$且$\{P\land\neg\,B\}\,S_2^*\,\{Q\}
    \vdash_{\mathrm{TO}}[P\land\neg\,B]\,S_2\,[Q]$, 则
    \vspace{1pt}
    \[\begin{aligned}
        &\,\,\{P\}\,\mathbf{if}\,\,B\,\,\mathbf{then}\,\,
        \{P\land B\}\,S_1^*\,\{Q\}\,\,\mathbf{else}\,\,\{P\land\neg\,B\}\,S_2^*
        \,\{Q\}\,\,\mathbf{fi}\,\,\{Q\}\\[-3pt]
        \,&\vdash_{\mathrm{TO}}[P]\,\mathbf{if}\,\,B\,\,\mathbf{then}\,\,S_1
        \,\,\mathbf{else}\,\,S_2\,\,\mathbf{fi}\,[Q],\\[2pt]
    \end{aligned}\]
    \item 如果$P\to I$, $\{I\land B\}\,S^*\,\{I\}
    \vdash_{\mathrm{TO}}[I\land B]\,S\,[I]$, $t,z$是整型变量,
    $z$不出现在$t,B,I,S$中, $\{I\land B\land t=z\}\,S^{**}\,\{t<z\}
    \vdash_{\mathrm{TO}}[I\land B\land t=z]\,S\,[t<z]$,
    $I\to t\geqslant 0$且$I\land\neg\,B\to Q$, 则
    \vspace{1pt}
    \[\begin{aligned}
        &\,\,\{P\}\{\mathbf{inv}:I\}\{\mathbf{bd}:t\}\,\mathbf{while}\,\,B\,\,
        \mathbf{do}\,\,\{I\land B\}\,S^*\,\{I\}
        \,\,\mathbf{od}\,\{Q\}\,\\[-3pt]
        \,&\vdash_{\mathrm{TO}}\,[P]\,\mathbf{while}\,\,B
        \,\,\mathbf{do}\,\,S\,\,\mathbf{od}\,[Q].
    \end{aligned}\]
\end{enumerate}

\vspace{15pt}

\hypertarget{sec:定理2.1}{}
\noindent\textbf{定理2.1 (正规形定理)}

对任意的断言$P,Q$和程序$S$, $\vdash_{\mathrm{TW}}[P]\,S\,[Q]$
当且仅当存在$\{P\}\,S^*\,\{Q\}\vdash_{\mathrm{TO}}[P]\,S\,[Q]$.

\vspace{15pt}

完全正确性注释程序上的操作几乎与定义1.1相同(只是多了$\mathbf{bd}$标记),
定理1.2的四条性质对完全正确性注释程序也依旧成立.

\end{document}
